\section{Torus charts and phasewise localization}\label{sec:torus-routing}

The next constructions keep the whole motion inside a specified torus
chart. A scaled linear interpolation preserves a convex first coordinate;
an exponential interpolation in a fixed chart instead uses endpoint-width
bounds. A phasewise construction then exhibits the vertical and planar
potentials separately, and the final variant retains a permanent origin
track. These choices give different path estimates for their observers.
In every case the force is recomputed at the prescribed physical viscosity
after any spatial rescaling.

\subsection{Physical chart scaling and thickened torus boxes}

\begin{proposition}[Linear scaling into the torus]
\label{prop:bb-mid-scaled-torus}
The direction recorder without its passive track has a unit-torus
realization periodic after $t=1$, with fixed label $(1/8,0,0)$ and
halting event $x_1\in(-1/4,0)\pmod1$.  Each period acts on full
solid boxes; the pressure is zero and the force has zero mean.
\end{proposition}
\begin{proof}\label{proof:bb-mid-scaled-torus}
Use the controls of Lemma~\ref{lem:bb-direction}. Let
$\Sigma=\Gamma\times(\{\bot,F\}\sqcup\mathcal R)\times\{0,1\}$
be its full cell alphabet after deleting the singleton passive track.
Let $S$ be the finite set of recorder controls and choose a bijection
$j:S\to\{1,\ldots,|S|\}$.
Use odd digits with $B=2|\Sigma|+1$, and offsets $o_s=-2j(s)$ for halting main
controls and $o_s=2j(s)$ for the others. Number the resulting reciprocal branches
$i=1,\ldots,N$, and write $k_i$ for branch $i$'s first-coordinate scale.
Thicken each reciprocal branch by
$[-1,1]$.  Its horizontal family gaps are at least $B^{-2}$.
Lift branch $i$ to $4i$, use the linear first scale
$\kappa_i=1+\theta(k_i-1)$, and descend.  Padding smaller than
$(4B^2)^{-1}$ separates all phases: in the middle the half-height
is one and height gaps are four; in the other phases horizontal
gaps suffice.  The two horizontal half-widths are at most $1/2$.
Use the true-plateau cutoff of Lemma~\ref{lem:bb-curl}; it supplies
at least the full-box guarantee, together with a neighborhood extension.
The exact first coordinate obeys
\begin{equation}\label{eq:bb-mid-scaled-convex}
 X_1(\tau)=(1-\theta)X_1(0)+\theta X_1(1).
\end{equation}

Let $V$ be the period-one extension of the unit-time field just constructed.
All codes, paths and padded supports lie in $(-L,L)^3$ for
$L=10+2|S|+4N$.  Put $a=1/(16L)$ and define in the torus patch
$(-1/2,1/2)^3$ the repeated field $U(t,x)=aV(t,x/a)$.
Its support lies in $(-1/16,1/16)^3$ and its curl potential is the
scaled potential $a^2A(t,x/a)$.  In particular its mean is zero.
The residual must be recomputed at the physical viscosity:
\begin{equation}\label{eq:bb-mid-viscosity-scale}
 \mathcal F_\nu[U](t,x)=a\{V_t+(V\cdot\nabla)V-(\nu/a^2)\Delta V\}(t,x/a).
\end{equation}
This keeps the specified viscosity rather than silently rescaling it.
During the first unit interval load from $(1/8,0,0)$ to the scaled
initial code using a curl translation with a box of half-width $1/32$
and padding $1/32$.  Its support lies inside $(-1/4,1/4)^3$.

For a nonhalting run both loading endpoints have positive first
coordinate, and every subsequent pair of coded endpoints is positive.
Equation~\eqref{eq:bb-mid-scaled-convex} preserves that sign throughout.
A halting checkpoint has negative first coordinate in $(-1/16,0)$.
All paths stay in the stated patch, so passage modulo one introduces
no extra visit.  Checkpoint times are
$1+\sum_{l<n}(5+2(l-h_l))$.  This proves the all-time equivalence,
including initial halting; Proposition~\ref{prop:bb-realization}
finishes the proof.
\end{proof}

\begin{proposition}[Exponential deformation in a fixed chart]
\label{prop:bb-mid-direct-torus}
The move-first recorder has a unit-torus realization periodic after
one loading interval, with fixed label $(1/4,1/4,1/4)$ and halting
event $x_1\in(2/3,5/6)\pmod1$.  The repeated motion acts on thickened
closed rectangles and their neighborhoods.  Its mean-zero residual
has the form $f_0+\nu f_1$ with coefficients independent of $\nu$.
\end{proposition}
\begin{proof}\label{proof:bb-mid-direct-torus}
First normalize to one halting work state and use the move-first recorder
of Lemma~\ref{lem:bb-movefirst}; let $\Sigma$ be its full three-track
cell alphabet. Initial halting is preserved by this normalization.
For $m$ controls take $\ell=1/[64(m+1)]$.  Nonhalt squares have
lower left corners $(1/4+2j\ell,1/4)$, $0\le j\le m-2$;
the halt square has lower left corner $(3/4,1/4)$.
Use odd digits and $b=2|\Sigma|+1$, with coding height $1/4$.
There are $N$ reciprocal branches with gaps at least $\ell b^{-2}$.
Thicken them by half-height $h_0=1/[64(N+1)]$ and choose heights
$z_i=1/2+i/[4(N+1)]$.  Use exponential scale $k_i^\theta$,
center interpolation, and
$\epsilon=\min(h_0/2,\ell/(8b^2))$.
The first and last phases have horizontal gaps exceeding
$2\epsilon$; middle height gaps equal $16h_0$ and exceed
$2(h_0+\epsilon)$.  Widths never exceed their endpoint values.
A concrete plateau cutoff for a box centered at $C_i$, with
half-widths $w_{ij}$, is
\[
 \chi_i=\prod_{j=1}^3\sigma\!\left(
 \frac{(w_{ij}+\epsilon)^2-(x_j-C_{ij})^2}
 {(w_{ij}+\epsilon)^2-(w_{ij}+\epsilon/2)^2}\right).
\]
Use its product with
$\tfrac12\dot C_i\times(x-C_i)+(0,0,(\dot a_i/a_i)(x_1-C_{i1})(x_2-C_{i2}))$,
where $a_i=k_i^\theta$.  Its curl gives the exact affine motion on
an open neighborhood of each box.  All supports remain inside the
unit cube.

Load the initial rational code along its segment from the fixed
label, using a radial cutoff of radius $1/32$ and a flat switch
$\sigma(4t-1)$.  Then repeat the unit motion, with phase switches
$\sigma(8\tau-1),\sigma(8\tau-3),\sigma(8\tau-5)$.
These choices give spatial zero collars of width $1/8$ and temporal
zero collars of width $1/16$ at nonnegative integer joins, including
loading.  The force follows from
$f_0=U_t+(U\cdot\nabla)U$, $f_1=-\Delta U$.

Nonhalt centers have first coordinate at most $1/4+1/32$ and
half-width at most $\ell/2$, hence their entire trajectories satisfy
$x_1\le1/4+1/32+\ell/2<2/3$.  This is a width bound, not a convexity
claim for exponential scaling.  Loading also avoids the slab for a
nonhalting input.  Halt codes enter it at a finite integer sample,
including time one for an initial halt.  The record block length and
the known frontier site $n+2$ recover absolute head position at each
work checkpoint.  All force and uniqueness assertions now follow
from Proposition~\ref{prop:bb-realization}.
\end{proof}

\subsection{Thin-sheet curls and a slab observer}
\label{ba:thin-sheet}

The next realization stays directly inside the unit cube.  It uses a
vertical translation potential in the first and last phases and a planar
Hamiltonian potential in the middle.  Because its observer ignores
height, the proof must keep the first coordinate safe throughout all
three phases.

\begin{theorem}[A direct unit-cube implementation]
\label{ba:thin-sheet-torus}
At each fixed positive computable viscosity, every machine and finite
input has an effective smooth mean-zero force on the unit flat
three-torus, with bounded mixed derivatives and period one for $t\ge1$,
whose smooth zero-data solution satisfies
\[
 \exists t\ge0:\ X(t;(1/4,1/4,1/4))\in
 (2/3,5/6)\times\mathbb T^2
 \quad\Longleftrightarrow\quad\text{the machine halts}.
\]
The normalized pressure is zero and the solution is unique in the
classical torus comparison class of Section~\ref{sec:model}.
\end{theorem}

\begin{proof}\label{ba:thin-sheet-torus-proof}
For the cutoffs in this proof, recall
\[
 E(r)=\begin{cases}e^{-1/r},&r>0,\\0,&r\le0,\end{cases}
 \qquad \Theta(r)=\frac{E(r)}{E(r)+E(1-r)}.
\]
Use the single-halt normalization and the seven-row recorder in
Lemma~\ref{lem:bb-movefirst}, with frontier initially at two and all marks
initially zero.  In the present notation its controls are
$S_q,U_r,W_r,C_q,D_q$ (ready, work-move completed, right scan,
frontier extension, left scan), and its full symbols are $(a,g,m)$
with history $g\in\{E,F\}\sqcup\{\widehat r\}$.
The first row writes and moves, the second marks and starts the scan,
and the last clears the mark and stays.  These are exactly the seven
rows already proved, after renaming.  Thus the new work head $j'$ is
less than $n+2$, a work step takes $2(n+2-j')+4$ local steps, and
both full-domain incoming properties hold.

Let $M$ count the partial controls, let $H=S_h$, and put
\[
 \ell=\frac1{100(M+1)},\qquad o_H=(3/4,1/2),\qquad
 o_{s_i}=(1/4+2i\ell,1/2)\quad(0\le i<M-1),\qquad z_*=1/4.
\]
In the state frames $o_s+\ell[0,1]^2$, use odd digits in base
$B=2|\Sigma|+1$ for the full alphabet $\Sigma$.  The minimal
left-split subdivision in Lemma~\ref{ba:rectangles} gives separately
separated rectangles $E_i^-,E_i^+$, centers $c_i^-,c_i^+$, and maps
\[
 y\longmapsto c_i^+
 +\operatorname{diag}(a_i,a_i^{-1})(y-c_i^-),
 \qquad a_i\in\{B^{-1},1,B\}.
\]
The common frame scale leaves these linear parts unchanged; all
rectangle side lengths are at most $\ell$.

For $N$ branches set $z_i=1/2+i/(4(N+1))$, $1\le i\le N$.
Let $\delta$ be one tenth of the minimum of $1/100$, all pairwise
source gaps, all pairwise target gaps and all $|z_i-z_j|$; omit empty
pair lists.  For a box $D=\prod_j[l_j,r_j]$, allowing zero widths,
define
\[
 \chi_D^\delta(x)=\prod_{j=1}^3
 \Theta\!\left(2(x_j-l_j+\delta)/\delta\right)
 \Theta\!\left(2(r_j+\delta-x_j)/\delta\right).
\]
It equals one on the $\delta/2$ padding and is supported in the
$\delta$ padding.  Put $\theta(s)=\Theta(3s-1)$.
On the first third of a unit step use
\[
 z_i(s)=z_*+(z_i-z_*)\theta(3s),\qquad
 D_i(s)=E_i^-\times\{z_i(s)\},\qquad
 P_i=\chi_{D_i(s)}^\delta(0,\dot z_i(s)x_1,0).
\]
Its curl on the plateau is $(0,0,\dot z_i)$; source gaps separate
these supports.  Hence every point of the whole source plate lifts
with planar coordinates unchanged.

On the middle third set
\[
 \eta=\theta(3s-1),\quad c_i=(1-\eta)c_i^-+\eta c_i^+,
 \quad\mu_i=1-\eta+\eta a_i,
\]
\[
 D_i(s)=
 \big[c_i+\operatorname{diag}(\mu_i,\mu_i^{-1})(E_i^--c_i^-)\big]
 \times\{z_i\}.
\]
With $\gamma_i=\dot\mu_i/\mu_i$, use
\[
 P_i=\chi_{D_i(s)}^\delta(0,0,\Psi_i),\qquad
 \Psi_i=\dot c_{i1}(x_2-c_{i2})-\dot c_{i2}(x_1-c_{i1})
          +\gamma_i(x_1-c_{i1})(x_2-c_{i2}).
\]
Its plateau curl is
$(\dot c_{i1}+\gamma_i(x_1-c_{i1}),
  \dot c_{i2}-\gamma_i(x_2-c_{i2}),0)$,
which is precisely the reciprocal-scale affine velocity.  The distinct
heights separate the supports during this phase.  In the last third use
the vertical potential again, with $E_i^+$ and
$z_i(s)=z_i+(z_*-z_i)\theta(3s-2)$; target gaps separate the descents.
Summing the curls gives the required field in each phase, with zero
collars at every join.

All supports lie strictly inside the unit cube: horizontal centers lie
in $[1/4,3/4+\ell]\times[1/2,1/2+\ell]$, half-widths are at most
$\ell/2$, heights lie between $1/4$ and $3/4$, and
$\delta\le10^{-3}$.  Thus the field extends smoothly and with zero
mean to the torus.  The construction acts on neighborhoods, not only
plates.  An initial perturbation of max norm less than $\delta/(4B)$
stays within $\delta/4$ of its affine plate path: vertical phases
preserve the perturbation, and both horizontal scale factors are at
most $B$.  It remains on the plateau, so uniqueness of the smooth ODE
gives that same affine continuation throughout.

Write the rational initial code as $(b,z_*)$.  During one unit of time
load the fixed particle using
\[
 c_0(t)=(1-\theta(t))(1/4,1/4)+\theta(t)b
\]
and the curl of
\[
 \chi_{\{(c_0(t),z_*)\}}^{1/100}
 (0,0,\dot c_{01}(x_2-c_{02})-\dot c_{02}(x_1-c_{01})).
\]
This has plateau velocity $(\dot c_0,0)$, support inside the cube,
and zero endpoint collars.  Repeat the step field thereafter and define
its residual force at $\nu$.  The finite compact phase formulas give
effective mixed-derivative bounds; the only scale denominators satisfy
$\mu_i\ge\min(1,a_i)>0$.  Curls have zero mean, and incompressibility
gives zero mean for the residual.  Lemma~\ref{ba:realization} supplies
the unique zero-data PDE solution and its global trajectories.

At times $1+k$, induction through the full branch maps gives the exact
local code through the first halt, or for all $k$ in the nonhalting
case.  A halt code has first coordinate in
$[3/4,3/4+\ell]\subset(2/3,5/6)$, including a code reached solely by
loading an initially halted machine.  For a nonhalting input the loader
stays below $1/2$ in first coordinate.  Lift and descent retain the
nonhalting planar endpoint codes.  In the middle phase the centers are
convex combinations of nonhalting centers and the first half-width is
at most $\ell/2$, so
\[
 X_1\le1/4+2M\ell+\ell/2<1/2.
\]
This excludes the height-independent observer at every intermediate
time.  Only the loader depends on the input; the repeated field depends
on the finite machine table.
\end{proof}


\subsection{A permanent origin marker and moving cutoff neighborhoods}
\label{bc:height-separated-section}
\begin{theorem}[A height-separated mechanism with an origin track]
\label{bc:height-separated-event}
On the unit torus, the fixed label $(1/8,3/8,1/4)$ and strip
$2/3<x_1<7/8$ realize halting by a smooth mean-zero general force,
periodic after time one, with zero initial velocity and pressure and bounded
mixed derivatives. A permanent tape track identifies the original absolute
origin. For a fixed machine the repeated field is independent of the input.
\end{theorem}
\begin{proof}\label{bc:height-separated-event-proof}
Normalize the work machine to one halting state as in
Section~\ref{sec:bb-frontiers}, preserving initial halting.
Use four tracks: work, head mark, history, and an origin bit. The history
is empty except for a frontier at cell two; the origin bit is one only
at absolute cell zero. In a ready control $S_q$, the head mark is one
at the work head. At intermediate controls use the seven-row recorder
unchanged, with the origin track copied in every row. The first row now
requires ready mark one and writes mark zero at the old head before the
work move. The final return row leaves mark one when entering $S_q$.
These are the only two changes to the seven rows.

For precision, this table is the conjugate of the frontier table on its
full configuration domain by the following involution: toggle the mark
at the current head exactly when the control is ready, and leave it
unchanged at every other control. Also take the direct product with the
passive origin track. This is a bijection of all configurations, not merely
of initialized ones. It proves the full-domain inverse. The incoming-rule
properties also persist: into the work-tag control the first-row output
mark is zero; into the right scan the marking row and its loop still write
one and zero; into the return scan frontier and nonfrontier outputs remain
distinct; into the ready control only the last row applies. The initial
configuration is the conjugate of the frontier initialization. Consequently
its cycle length is $2(2+n-h')+4$, its absolute frontier is $2+n$, and
its origin bit remains fixed at the absolute origin. This verifies the
distinct initialization without another appeal to an existential compiler.

Denote the resulting full four-track alphabet by $\mathcal A$, and let
$K=2|\mathcal A|+1$ be its odd-digit radix.
For $m$ controls take state-square side $r=1/[16(m+1)]$, with nonterminal
lower corners $(1/8+2(i-1)r,1/8)$ and terminal lower corner $(3/4,1/8)$.
Their common coding height is $z_0=1/4$. Lemma~\ref{lem:bb-onecell} supplies
separated full planar rectangles; let $N$ be the number of branches.
Give branch $i$ height
$h_i=1/2+i/[4(N+1)]$. During a first phase lift its center from $z_0$
to $h_i$, keeping its planar shape; during a middle phase interpolate
its planar center, with width
$w_i=(1-\mu)w_i^-+\mu w_i^+$ and height
$v_i=w_i^-v_i^-/w_i$; during a final phase lower it to $z_0$.
The planar first coordinate is again an endpoint convex combination.
Choose a fixed positive padding less than one tenth of all source and
target gaps, all chart-face margins and all height separations. The
moving padded rectangles are disjoint at every phase. During the middle
phase this follows from the different heights; during the first and last
phases it follows from endpoint planar separation.

Here is a direct potential for these moving neighborhoods. With moving
center $c=(c_x,c_y,c_z)$ and $\lambda=\dot w/w=-\dot v/v$, set
\[
 A=(0,\dot c_z(x-c_x),
       \dot c_x(y-c_y)-\dot c_y(x-c_x)
                         +\lambda(x-c_x)(y-c_y)).
\]
Its curl on the plateau is
$(\dot c_x+\lambda(x-c_x),\dot c_y-\lambda(y-c_y),\dot c_z)$.
Multiply it by a cutoff supported in the moving padded neighborhood and
sum the curls. All supports remain in the open cube. The displayed
center-and-shape paths solve this field on the full planar rectangles;
uniqueness proves their exact endpoint maps. The positive padding and
bounded planar scales also give a uniform initial neighborhood that stays
on every plateau, by the perturbation argument of Lemma~\ref{lem:bb-curl}.
Each component of the total
field has zero mean after periodization, since the entire field is a curl.

For initialization, choose equal small squares about the fixed label and
the rational input code at height $z_0$, with half-side less than one
quarter of their minimum chart-face margin. Move one to the other by the
same localized potential for translation during $[0,1]$. A nonterminal
loading first coordinate lies between $1/8$ and $1/4$, and every subsequent
nonterminal phase lies below $1/4$. A terminal endpoint lies in the detector.
The initialized symbolic equivalence, including initial halting, therefore
holds at all real times. Repeating the finite middle mechanism and applying
Lemma~\ref{bc:analytic} gives the asserted force and comparison properties.
Only the loading translation and its input code depend on the finite word.
\end{proof}
